% Part: model-theory
% Chapter: models-of-arithmetic
% Section: introduction

\documentclass[../../../include/open-logic-section]{subfiles}

\begin{document}

\olfileid{mod}{mar}{int}
\section{Pambuka}

\emph{Model standar} aritmetika
iku !!{structure}~$\Struct{N}$
kanthi $\Domain{N} = \Nat$,
ing ngendi $\Obj{0}$,
$\prime$, $+$, $\times$,
lan $<$ diinterpretasi
kaya kang lumrah.
Tegesé, $\Obj{0}$ iku $0$,
$\prime$ iku fungsi penerus,
$+$ diinterpretasi minangka
panjumlahan lan $\times$
minangka ping-pingan
wilangan ing~$\Nat$.
Luwih cetha,
\begin{align*}
  \Assign{\Obj{0}}{N} & = 0\\
  \Assign{\prime}{N}(n) & = n + 1\\
  \Assign{+}{N}(n, m) & = n + m\\
  \Assign{\times}{N}(n, m) & = nm
\end{align*}
Mesthi ana struktur-struktur
kanggo $\Lang{L_A}$ kang
ranahé dudu~$\Nat$.
Upamane, kita bisa njupuk
$\Struct{M}$ kanthi ranah
$\Domain{M} = \{a\}^*$
(urutan winates saka siji
simbol~$a$, yaiku
$\emptyset$, $a$, $aa$,
$aaa$, \dots), lan
interpretasi iki:
\begin{align*}
  \Assign{\Obj{0}}{M} & = \emptyset\\
  \Assign{\prime}{M}(s) & = s \concat a\\
  \Assign{+}{M}(a^n, a^m) & = a^{n + m}\\
  \Assign{\times}{M}(a^n, a^m) & = a^{nm}
\end{align*}
Cathetan koreksi sumber OLPL-094:
rong operasi pungkasan nampa
urutan saka ranah struktur,
dene eksponen nyatakake
dawane urutan; sumber nulis
inpute minangka angka.
Rong struktur iki sajatiné
padha: bedane mung
!!{element}s saka
!!{domain}s-e, dudu
carane !!{element}s
saka !!{domain}s iku
sesambungan siji lan
sijiné lumantar
fungsi interpretasi.
Kita ngarani rong
!!{structure}s iki
\emph{isomorfis}.

Akibat gampang saka teorema
kompaktitas yaiku yen sembarang
teori kang bener ing~$\Struct{N}$
uga duwe model kang ora isomorfis
karo~$\Struct{N}$.
Struktur kaya mangkono
diarani \emph{nonstandar}.
Kang narik kawigaten,
!!{element}s saka model standar
(yaiku $\Struct{N}$, lan uga
kabeh !!{structure}s kang
isomorfis karo iku)
kabeh dadi nilai angka
standar~$\num{n}$, yaiku
\[
\Domain{N} = \Setabs{\Value{\num{n}}{N}}{n \in \Nat}
\]
Nanging ora kaya mangkono
ing model nonstandar:
yen $\Struct{M}$ nonstandar,
mula ana paling ora siji
$x \in \Domain{M}$ saéngga
$x \neq \Value{\num{n}}{M}$
kanggo saben~$n$.

Unsur-unsur nonstandar iki
narik kawigaten: bisa dianggep
minangka wilangan asli
tanpa wates. Nanging
anane uga nerangake,
ing sawijine teges,
gejala ora-komplet.
Gatekna contone pratelan
konsistensi aritmetika Peano,
$\OCon[\Th{PA}]$, yaiku
$\lnot
\lexists[x][\OPrf[\Th{PA}](x, \gn{\lfalse})]$.
Amarga $\Th{PA}$ ora
mbuktekake $\OCon[\Th{PA}]$
utawa $\lnot \OCon[\Th{PA}]$,
salah siji saka loro iku
bisa ditambahake marang
$\Th{PA}$ tanpa
ngrusak konsistensi.
Amarga $\Th{PA}$
konsisten,
$\Sat{N}{\OCon[\Th{PA}]}$,
lan mula
$\Sat/{N}{\lnot
  \OCon[\Th{PA}]}$.
Dadi $\Struct{N}$
\emph{dudu} model saka
$\Th{PA}
\cup \{\lnot \OCon[\Th{PA}]\}$,
lan kabeh modele
kudu nonstandar.
Model saka
$\Th{PA} \cup
\{\lnot \OCon[\Th{PA}]\}$
kudu ngemot sawenehing
!!{element} kang dadi
saksi supaya
$\lexists[x][\OPrf[\Th{PA}](x, \gn{\lfalse})]$
bener, yaiku nomer
G\"odel saka !!a{derivation}
kontradiksi saka~$\Th{PA}$.
Cathetan koreksi sumber OLPL-095:
predikat bukti ing ukara
eksistensial iki kudu
ngemot argumen saksi kang
wis diiket kuantor.
!!a{element} kaya mangkono
ora bisa standar---amarga
$\Th{PA} \Proves \lnot
\OPrf[\Th{PA}](\num{n}, \gn{\lfalse})$
kanggo saben~$n$.

\end{document}
